\subsection{Sensitivity of the measured ancestry-response test}
\label{note:measured_alignment_power}
We estimated detection sensitivity conditional on the observed sampling, using a Gaussian response model frozen before simulation. The analysis retained 125 partner observations of 102 presynaptic roots across thirteen scans and seven targets, including twenty target--partner combinations repeated across scans and two roots shared between targets. Each scan contained 136 repeated stimuli: 130 presented twice and six ten times. Original condition hashes identified 316 distinct repeated stimuli across scans. Pair tables, contact geometry, recording identities and the required observed responses are included in Source Data.

For each target arbor, an ancestry covariance was defined by the shared soma-to-common-ancestor path length divided by the geometric mean of the two soma-to-contact lengths. This is a positive-semidefinite Brownian path kernel; zero-depth contacts contributed independent diagonal components. The seven kernels were embedded in common presynaptic-root coordinates, summed and diagonal-normalized, so each shared root had one consistent latent response. Reliable tuning mixed independent Gaussian unit-specific and ancestry-correlated components with variance fractions $1-\lambda$ and $\lambda$, respectively. We tested $\lambda=0,0.05,\ldots,1$, retaining the actual partner sets and stimulus identities.

Measurement noise was calibrated to each recording's odd/even split-half Spearman reliability. Values were clipped to $[0,1]$ and converted to the Gaussian Pearson parameter $r_p=2\sin(\pi r_s/6)$; the two negative estimates remained archived and contributed zero reliable variance. With $h$ the mean inverse half-repeat count, trial responses had reliable variance $r_p$ and noise variance $(1-r_p)/h$. Condition-mean noise was divided by the actual repeat count. This matches the pooled Pearson second moment; its Spearman calibration is approximate because repeat counts vary. An independent 1,000-dataset audit found mean and maximum absolute Spearman discrepancies of 0.0020 and 0.0069, without retuning. A prespecified perfect-reliability comparison removed measurement noise.

In 4,000 simulated complete datasets, all pairwise similarities were computed from joint response matrices. The original partial-rank statistic controlled ranked Euclidean separation and contact-depth difference; scan effects were averaged within target. Positive detection required a positive mean effect and an exact two-sided signed-rank $P\leq0.05$ over the seven targets. All 128 sign assignments were enumerated. For seven nonzero untied effects, the minimum two-sided $P$ is 0.015625, and six of 128 null sign assignments reject at 0.05. This discreteness does not imply a universal power ceiling below one. Gaussian draws were paired across signal fractions and reliability scenarios; Wilson intervals describe Monte Carlo uncertainty.

With measured reliability, the first sustained tested 80\% crossing was $\lambda=0.60$, bracketed by 0.55--0.60. Descriptive interpolation gave $\lambda=0.589$ and a mean simulated target-level partial-rank effect of 0.249. Detection at $\lambda=0.60$ was 80.85\% (95\% Monte Carlo interval, 79.60--82.04\%); the lower interval first exceeded 80\% at $\lambda=0.65$. At maximal specified signal, detection was 96.23\% (95.59--96.77\%). The zero-signal two-sided rejection rate was 4.95\% (4.32--5.67\%). Perfect reliability reduced the interpolated variance fraction to 0.206, with a similar mean partial-rank effect of 0.248 at detection. These sensitivity estimates assume the specified Gaussian ancestry covariance, reliability calibration and known shared-input dependence; they do not exclude weaker effects, other covariance patterns, unknown biological dependence or alignment involving unobserved partners.
